\begin{quadro}[htbp]
\centering\small
\caption{Termos da organização experimental}\label{qua:termos}
\begin{tabular}{|p{0.18\textwidth}|p{0.71\textwidth}|}
\hline
Termo & Significado neste trabalho\\ \hline
Evento & Um registro sintético enviado à API. Possui identificador único; milhares de eventos podem pertencer à mesma coleta.\\ \hline
Coleta & Uma execução de uma variante em uma condição, com aquecimento, medição e observação posterior. É a unidade experimental, não um evento.\\ \hline
Condição & Combinação de perfil de carga ou falha, taxa e variante. Foram treze condições com cinco coletas cada.\\ \hline
Par & Duas coletas, uma por variante, com o mesmo conjunto determinístico de eventos na mesma condição comparável.\\ \hline
Bloco & Agrupamento de uma repetição de todas as condições, usado para organizar e contrabalançar a ordem. Foram cinco blocos.\\ \hline
Piloto & Verificação prévia dos instrumentos e parâmetros. Não integra as repetições definitivas.\\ \hline
Revisão analítica & Emenda posterior ao início da campanha que definiu elegibilidade por indicador diante de uma evidência temporal incompleta, sem alterar os dados ou repetir a coleta.\\ \hline
\end{tabular}
\fonte{Elaborado pelo autor (2026).}
\end{quadro}
